% ============================================================
% JPDC extension --- manuscript skeleton (zero-GPU draft).
%
% Status: planning skeleton only.  Every W1 result is a placeholder.
% No W1 outcome, trace value, or contrast has been inserted.
%
% Overlap discipline: see overlap_ledger.md.  The signed Springer
% proceedings licence (CP025/CP025_COPYRIGHT.pdf, clause 4(d)(v))
% permits an expanded version with at least 30% new material and an
% express statement of the incremental change.  The target budget in
% the ledger is deliberately above that floor.
%
% Build (once elsarticle.cls is available):
%   pdflatex jpdc_skeleton
%   bibtex   jpdc_skeleton
%   pdflatex jpdc_skeleton
%   pdflatex jpdc_skeleton
% ============================================================
\documentclass[preprint,12pt]{elsarticle}

\usepackage[T1]{fontenc}
\usepackage[utf8]{inputenc}
\usepackage{amsmath,amssymb}
\usepackage{booktabs}
\usepackage{graphicx}
\usepackage{siunitx}
\usepackage{url}
\usepackage{xcolor}
\usepackage{microtype}
\usepackage{placeins}

% Placeholder used for every W1 result.  Keep it visually obvious in drafts.
\newcommand{\wone}[1]{\textcolor{red}{[W1 PLACEHOLDER: #1]}}
\newcommand{\ppam}{the PPAM~2026 paper}
\newcommand{\mf}{\mathrm{MF}}
\newcommand{\bp}{\mathrm{BP}}

\journal{Journal of Parallel and Distributed Computing}

\begin{document}

\begin{frontmatter}

\title{From Per-Pass Savings to End-to-End Energy:\\
A Pre-Registered Study of Forward-Only Learning}

\author[a]{Przemys{\l}aw Spyra\corref{cor1}}
\ead{pspyra@agh.edu.pl}
\author[a]{Witold Dzwinel}
\affiliation[a]{AGH University of Krakow, Faculty of Computer Science,
al.~Mickiewicza 30, 30-059 Krak\'ow, Poland}
\cortext[cor1]{Corresponding author.}

\begin{abstract}
Backpropagation-free training is often described as energy-efficient because
it removes the backward sweep, but a per-pass saving does not imply an
end-to-end saving: a method that needs more passes can consume more energy
overall.  This paper separates the two effects.  We define the per-pass energy
advantage $a$ and the pass multiplier $m$, derive the condition $a>m$ under
which a forward-only method wins end to end, and test that condition in a
pre-registered, hardware-instrumented family of experiments on NVIDIA A100
GPUs.  The design covers five training arms, two image-classification
datasets, twenty paired seeds, tuned BP controls, a scheduled-BP control, and
budget-matched controls.  \wone{insert the W1 end-to-end result and its
pre-registered uncertainty statement here; do not insert a number before the
frozen analysis is complete.}  The contribution is a measurement framework and
a confirmatory test, not a FLOP-count argument.
\end{abstract}

\begin{keyword}
forward-only learning \sep Mono-Forward \sep backpropagation \sep
energy measurement \sep pre-registration \sep reproducibility
\end{keyword}

\end{frontmatter}

% ------------------------------------------------------------------
% Section map (details and word budget: overlap_ledger.md)
%  1 Introduction                         revised
%  2 Background and related work          revised
%  3 End-to-end energy theory             NEW
%  4 Mono-Forward and the cost hypothesis revised
%  5 Measuring cost fairly                revised
%  6 W1 pre-registered design             NEW
%  7 Results                              mostly NEW
%  8 Discussion and limitations           revised
%  9 Conclusion                           revised
%  Appendices                             NEW
% ------------------------------------------------------------------

\section{Introduction}
\label{sec:intro}

Deep neural networks are trained overwhelmingly with backpropagation
(BP)~\cite{rumelhart1986learning}.  BP is effective, but it couples the layers: the
backward sweep must retain activations and must wait for the layer above.
Forward-only algorithms replace the global backward sweep with local learning
signals.  Hinton's Forward-Forward (FF)~\cite{hinton2022forward}, Cascaded Forward
(CaFo)~\cite{zhao2023cafo}, and Mono-Forward (MF)~\cite{gong2025mono} are the
three representatives considered here.  Their appeal is architectural as well
as numerical: no global backward graph, no activation store for a backward
sweep, and a natural mapping to pipelined or memory-constrained hardware.

The efficiency case is usually made with operation counts.  That is not
enough.  A method can execute fewer FLOPs and still consume more energy if it
runs longer, uses the device poorly, or moves more data.  The PPAM~2026 paper
measured energy on hardware and found that removing the backward pass is not by itself
a saving: FF and CaFo consumed more energy than a tuned BP baseline in the
tested configurations, while MF was the exception in the matched-data-pass
protocol.  The same study also found that MF needed more passes to reach its
own convergence, which is exactly the regime in which a per-pass advantage can
disappear end to end.

This paper extends that work in four ways.  First, it makes the end-to-end
argument explicit: Section~\ref{sec:theory} defines per-pass energy, pass
multiplier, and the win condition, including fixed-cost and uncertainty
corrections.  Second, it reports a pre-registered confirmatory family (W1) with
five arms, two datasets, twenty paired seeds, tuned BP controls, a scheduled-BP
control, and budget-matched controls; the protocol was frozen before any W1
outcome was analysed.  Third, it adds a measurement protocol that separates
matched-data-pass energy from matched-accuracy energy, so that a per-pass
result cannot be presented as an end-to-end result.  Fourth, it releases the
frozen configurations, analysis code, and table shells needed to reproduce the
comparison.

\textbf{Incremental contribution.}  This manuscript is an expanded version of
\ppam.  The conference paper established the hardware measurement and the
matched-data-pass comparison.  The journal version adds the formal cost model
(Section~\ref{sec:theory}), the pre-registered W1 design
(Section~\ref{sec:w1}), the end-to-end and control analyses
(Section~\ref{sec:results}), the budget-matched Protocol~B analysis, and the
extended discussion of validity.  The exact incremental statement will be
repeated in the cover letter at submission.

\section{Background and Related Work}
\label{sec:background}

\subsection{Backpropagation and its cost}
BP computes a global gradient by a forward sweep followed by a backward sweep.
The backward sweep requires the activation of each layer to remain available
until its gradient has been computed.  The resulting memory and synchronisation
costs are well known; the energy cost is less often measured.  \wone{cite the
specific BP energy literature used in the final related-work section; do not
reuse the PPAM bibliography without checking each entry.}

\subsection{Forward-only algorithms}
FF replaces the global objective with layer-local goodness objectives.
CaFo cascades local classifiers and trains them sequentially.  MF combines a
forward stage with an auxiliary objective and a detached path between stages.
The three methods are not interchangeable: they differ in architecture,
parameter count, local objective, and the number of forward stages per data
pass.  A fair comparison therefore has to fix the data-pass definition and the
tuning budget, not only the nominal FLOP count.

\subsection{Energy measurement}
Hardware energy is the integral of instantaneous power over time.  NVML
provides device-level power samples at a fixed interval; the integral is an
estimate whose error depends on the sampling interval, the clock behaviour, and
the treatment of idle periods.  The W1 traces retain the raw samples and the
segment boundaries so that the estimator can be audited.  \wone{add the
measurement-error discussion and the exact NVML version used in the final
runs.}

\subsection{Fair tuning and pre-registration}
A comparison between a tuned baseline and a default-configured alternative is
not a comparison of algorithms.  The W1 design freezes the search space, the
trial budget, the sampler, and the stopping rule before the confirmatory runs.
The BP-bestval control gives BP the best validation-selected configuration from
its own search; the BP-scheduled control adds a cosine schedule and disables
early stopping; the budget-matched controls equalise the pass budget using the
frozen Protocol~B rule.  These controls address different objections and are
reported separately.

\section{End-to-End Energy of Forward-Only Training}
\label{sec:theory}

% ============================================================
% JPDC extension --- theory section draft.
% Status labels used throughout:
%   [Definition]  a quantity or measurement boundary is defined;
%   [Derivation]  the statement follows algebraically from the definitions;
%   [Conjecture]  an empirical claim that the W1 design can test only in part.
% No W1 outcome value appears in this file.  All result slots are symbolic.
% ============================================================

\section{End-to-End Energy of Forward-Only Training}
\label{sec:theory}

This section separates the two quantities that are often conflated in
discussions of backpropagation-free learning: the energy of one training pass
and the energy needed to reach a specified accuracy.  The distinction is
elementary, but it is the difference between a per-pass efficiency claim and an
end-to-end efficiency claim.  The W1 family is designed to measure both
quantities under one instrumentation boundary.

\noindent\emph{Status labels.}  A statement labelled \textbf{[Definition]}
fixes a quantity or a measurement boundary; \textbf{[Derivation]} means the
statement follows algebraically from the definitions; \textbf{[Conjecture]}
means the statement is an empirical hypothesis that the W1 design can test only
in part.  Unlabelled sentences are connective exposition and make no
independent empirical claim.  No W1 outcome value appears in this section.

\subsection{Definitions and measurement boundary}

\textbf{[Definition]} Fix a training set $\mathcal{D}$, a model, an
optimisation algorithm, and an evaluation metric $A$ (here, test accuracy).
Let $A^\star$ be a target accuracy fixed before the comparison.  A
\emph{data pass} is one traversal of $\mathcal{D}$ through the training
computation that the algorithm actually performs.  For backpropagation (BP),
one data pass comprises the forward sweep, the backward sweep, and the
parameter update.  For a forward-only algorithm, one data pass comprises every
forward stage, local loss, and local update that the algorithm performs during
that traversal.  This definition deliberately counts the algorithm as
implemented, not an idealised forward FLOP count.

\textbf{[Definition]} Let $e_f$ and $e_b$ be the mean energy consumed per data
pass by the forward-only method and by BP, respectively, measured at the same
device, software stack, batch size, and instrumentation boundary.  Let
$N_f(A^\star)$ and $N_b(A^\star)$ be the number of data passes each method needs
to reach $A^\star$ for the first time.  When the target is unambiguous we write
$N_f$ and $N_b$.

\textbf{[Definition]} The \emph{per-pass energy advantage} of the forward-only
method is
\begin{equation}
  a \;=\; \frac{e_b}{e_f},
  \label{eq:advantage}
\end{equation}
and its \emph{pass multiplier} is
\begin{equation}
  m \;=\; \frac{N_f}{N_b}.
  \label{eq:multiplier}
\end{equation}
Both quantities are dimensionless.  $a>1$ means that one forward-only pass is
cheaper than one BP pass; $m>1$ means that the forward-only method needs more
passes to reach the same target.  Neither statement alone determines the
end-to-end comparison.

\textbf{[Definition]} Ignoring fixed start-up and evaluation costs for the
moment, the end-to-end training energies are
\begin{equation}
  E_f \;=\; N_f\,e_f,
  \qquad
  E_b \;=\; N_b\,e_b .
  \label{eq:endtoend}
\end{equation}
Equations~\eqref{eq:endtoend} are the formal content of the phrase
``end-to-end energy is per-pass energy times the number of passes''.  The
per-pass energies are not assumed equal, and the pass counts are not assumed
equal.

\subsection{The win condition}

\textbf{[Derivation]} From \eqref{eq:advantage}--\eqref{eq:endtoend}, the ratio
of end-to-end energies is
\begin{equation}
  R \;=\; \frac{E_f}{E_b}
      \;=\; \frac{N_f}{N_b}\,\frac{e_f}{e_b}
      \;=\; \frac{m}{a}.
  \label{eq:ratio}
\end{equation}
Therefore, in the no-overhead model,
\begin{equation}
  E_f < E_b
  \quad\Longleftrightarrow\quad
  a > m .
  \label{eq:win}
\end{equation}
A forward-only method with a per-pass advantage $a$ wins end to end exactly
when its pass multiplier is smaller than that advantage.  If $a=m$, the two
methods have equal training energy in this model; if $a<m$, the forward-only
method is cheaper per pass and more expensive end to end.

\textbf{[Derivation]} The result is not a statement about FLOPs.  The
quantities $e_f$ and $e_b$ are measured energies, so \eqref{eq:win}
already includes power draw, achieved clock, memory traffic, and any idle time
inside the pass.  A method can execute fewer arithmetic operations and still
have a smaller $a$ if it utilises the device poorly or moves more data.

\textbf{[Derivation]} Fixed costs change the threshold.  Let
\begin{equation}
  E_f = c_f + N_f e_f,
  \qquad
  E_b = c_b + N_b e_b,
  \label{eq:overhead}
\end{equation}
where $c_f$ and $c_b$ collect start-up, data-loading, evaluation, and other
run-level costs.  Substituting \eqref{eq:advantage}--\eqref{eq:multiplier} and
rearranging gives
\begin{equation}
  a \;>\; \frac{m}
  {1 + \dfrac{c_b-c_f}{N_b e_b}} .
  \label{eq:win-overhead}
\end{equation}
When $c_f=c_b$, \eqref{eq:win-overhead} reduces to \eqref{eq:win}.  A
forward-only method with larger fixed costs must clear a higher per-pass
advantage; a method with smaller fixed costs can tolerate a larger pass
multiplier.  The sign of $c_b-c_f$ therefore matters, and the W1 traces retain
the run boundary and the recorded evaluation timestamps; where a fixed cost
cannot be separated, it is bounded rather than assumed zero.

\textbf{[Definition]} The \emph{matched-accuracy end-to-end energy} of a method
is the cumulative energy consumed from the start of training until the first
pass at which the evaluation metric reaches $A^\star$.  This is the quantity
compared in \eqref{eq:endtoend}.  If a method never reaches $A^\star$, it has no
finite matched-accuracy energy and is reported as a censored run rather than
silently assigned the energy of its last pass.

\textbf{[Conjecture]} Removing the backward sweep is not by itself sufficient
for an end-to-end gain.  The empirically relevant question is whether the
measured pair $(a,m)$ satisfies \eqref{eq:win}; a method that is cheaper per
pass can still lose if its convergence is slow enough.  This conjecture is the
organising hypothesis of the W1 family, not a result of this section.

\subsection{What the W1 design can test}

The cost model is useful only if its quantities are measurable.  Table
\ref{tab:cost-assumptions} lists the assumptions behind
\eqref{eq:win}--\eqref{eq:win-overhead} and the part of the W1 design that can
test each one.  The table is deliberately explicit about what is direct
evidence, what is partial evidence, and what is outside the fixed W1
architecture.

\begin{table}[t]
\centering
\caption{Assumptions of the end-to-end cost model and their W1 tests.  ``Direct''
means the W1 instrumentation records the required quantity; ``partial'' means
the design constrains the assumption but does not vary it; ``not tested'' means
the assumption is a scope limitation.  No outcome value is shown.}
\label{tab:cost-assumptions}
\small
\begin{tabular}{@{}p{0.30\linewidth}p{0.20\linewidth}p{0.42\linewidth}@{}}
\toprule
Assumption & W1 evidence & Comment \\
\midrule
$e_f,e_b$ measured at one boundary
  & direct
  & NVML energy integration with the same start/stop rule for every arm. \\
Same target $A^\star$
  & direct
  & Pre-registered target and evaluation protocol; censored runs reported. \\
$N_f,N_b$ counted as implemented
  & direct
  & Epoch/pass counters; no FLOP-based substitution. \\
Per-pass energy stable over training
  & partial
  & Per-stage sums versus total energy; drift check across seeds. \\
Fixed costs estimated
  & partial
  & Run boundary and evaluation timestamps; fixed costs bounded where not separable. \\
Same data pipeline and batch size
  & direct (design)
  & Frozen configs; no per-arm preprocessing. \\
Same hardware generation
  & partial
  & A100 partition, but node-to-node variation is recorded, not eliminated. \\
Scaling to larger widths/depths
  & not tested
  & W1 fixes the $3\times2000$ MLP; scale is a separate conjecture. \\
\bottomrule
\end{tabular}
\end{table}

\textbf{[Definition]} The \emph{measured} per-pass advantage is
$\hat a=\bar e_b/\bar e_f$, where the bars denote the protocol-level means over
the pre-registered seeds.  The \emph{measured} pass multiplier is
$\hat m=\bar N_f/\bar N_b$ at the same target.  These are estimators of
\eqref{eq:advantage} and \eqref{eq:multiplier}, not definitions of new
quantities.  Their uncertainty must be propagated: a point estimate
$\hat a>\hat m$ is not by itself evidence that $a>m$ in the population of
seeds, because both quantities are measured with run-to-run variation.

\textbf{[Derivation]} If $\hat a$ and $\hat m$ are approximately independent
with relative standard errors $s_a$ and $s_m$, the relative standard error of
$\hat R=\hat m/\hat a$ is approximately
$\sqrt{s_a^2+s_m^2}$.  A decision rule based on \eqref{eq:win} therefore needs
a confidence interval for $R$, not only a comparison of two point estimates.
The W1 analysis reports the interval and the paired-seed structure; the theory
section does not assume that the interval excludes one.

\textbf{[Conjecture]} The pair $(a,m)$ may be architecture-dependent in a way
that the fixed W1 MLP cannot reveal.  In particular, $a$ may improve with width
because activation memory and parameter traffic grow differently in the two
methods, while $m$ may improve with a better local objective or a schedule that
stops earlier.  A scale-dependent crossover, if it exists, would require a
separate width/depth study; W1 can falsify the claim that backward removal is
sufficient at its fixed scale, but it cannot establish a general scaling law.

\subsection{Interpretation rules}

\textbf{[Definition]} We call a result \emph{end-to-end favourable} to the
forward-only method when the matched-accuracy energy ratio satisfies
$R<1$ with a pre-registered uncertainty rule.  We call it \emph{per-pass
favourable} when $a>1$ but make no end-to-end claim from that fact alone.  This
terminology prevents the common slide from ``cheaper per pass'' to ``cheaper
overall''.

\textbf{[Derivation]} Three boundary cases follow directly from
\eqref{eq:ratio}.  First, if $a\le 1$, the forward-only method cannot win end to
end at any $m\ge 1$ in the no-overhead model.  Second, if $a>1$ but
$m\ge a$, the per-pass advantage is exhausted by the extra passes.  Third, if
$m<1$ as well as $a>1$, the forward-only method wins for two independent
reasons; this is a stronger claim than either effect alone and should be
reported separately.

\textbf{[Conjecture]} Early stopping and schedule design act primarily on $m$,
not on $a$.  A fair comparison should therefore report both the stopping rule
and the realised pass count; comparing a forward-only method at its own
convergence with a BP baseline at a fixed epoch budget confounds the two.
The W1 BP-scheduled and budget-matched blocks are designed to separate these
effects, but the conjecture that schedule improvements generalise across
datasets remains to be tested.

\subsection{Summary of the testable model}

The model has one definitional identity, \eqref{eq:endtoend}, one algebraic
win condition, \eqref{eq:win}, and one overhead correction,
\eqref{eq:win-overhead}.  Its empirical content is the pair $(a,m)$ measured at
a fixed target.  The W1 design measures $a$ in the matched-data-pass protocol,
measures $m$ from the pass count to the pre-registered target, retains the
segments needed for $c_f$ and $c_b$, and records node and GPU identity so that
hardware variation is visible rather than absorbed into the treatment effect.
The design does not vary width, depth, or device generation; those remain
conjectures for a later study.


\section{Mono-Forward and the Cost Hypothesis}
\label{sec:mf}

\subsection{The Mono-Forward update}
MF trains a stack of stages with a local objective.  In the variant used here,
each stage produces a representation and a local prediction; the local loss is
applied at the stage, and the path between stages is detached for the purpose
of the local update.  The detach is not an implementation detail: it is the
mechanism that removes the global backward graph.  The W1 ablation includes a
variant without the M0 stage and variants with joint, simultaneous, and
sequential training of the stages, so that the cost of the mechanism can be
separated from the cost of the architecture.

\subsection{Why a per-pass advantage is plausible}
A forward-only pass avoids the backward sweep and the activation retention it
requires.  It may also run at a lower power draw because it has fewer
synchronisation points.  These effects can make $a>1$.  They do not guarantee
that $m\le a$: local objectives can converge more slowly, and a method that
needs more passes can lose end to end even with a substantial per-pass
advantage.  Section~\ref{sec:theory} states the condition; the W1 family
measures both sides of it.

\subsection{The cost hypothesis}
\textbf{[Conjecture]} The end-to-end advantage of MF is not a consequence of
removing the backward pass alone.  It is jointly determined by the per-pass
energy ratio and the pass multiplier at the target accuracy.  A convincing
demonstration must report both, with the target and the stopping rule fixed in
advance.  This is the hypothesis that the W1 design is built to test.

\section{Measuring Cost Fairly: Two Protocols}
\label{sec:protocols}

\subsection{The matched-data-pass protocol}
In the matched-data-pass protocol, every arm performs the same number of data
passes and the comparison asks for the energy per pass and the accuracy reached
at that budget.  This protocol isolates the cost of the update rule.  It is the
protocol used in \ppam{} and it is the correct protocol for the question
``which algorithm is cheaper for the same amount of training data?''  It is not
an end-to-end protocol, because the arms need not reach the same accuracy.

\subsection{The matched-accuracy protocol}
In the matched-accuracy protocol, every arm trains until it first reaches the
pre-registered target $A^\star$, and the comparison asks for cumulative energy
and passes to that point.  Runs that never reach $A^\star$ are censored, not
assigned the energy of their last pass.  This protocol answers the end-to-end
question and is the basis of the win condition in Section~\ref{sec:theory}.

\subsection{Instrumentation boundary}
Energy is measured at the device level with NVML, using the same start and stop
events for every arm.  The traces record node identity, GPU UUID, start time,
per-stage boundaries, and the evaluation segments.  The analysis integrates
power over the training interval and reports the per-pass and cumulative
values separately.  \wone{insert the final instrumentation-boundary paragraph
and the NVML sampling interval; the frozen protocol fixes the boundary, not the
hardware generation.}

\subsection{Tuning and stopping}
The tuning budget, search space, sampler, and pruner are fixed per arm.  The
confirmatory runs use the selected configuration; the BP-bestval control uses
BP's best validation-selected configuration rather than a default.  The
BP-scheduled control changes only the schedule and the early-stopping setting,
so that its difference from BP-bestval is interpretable.  The budget-matched
controls use the Protocol~B rule, whose per-stage budgets are computed from the
core BP runs after those runs complete; the rule is frozen before the controls
are run.

\section{The W1 Pre-Registered Design}
\label{sec:w1}

\subsection{Questions and family}
W1 asks whether MF's per-pass advantage survives an end-to-end comparison, and
whether any advantage is robust to tuned BP, a scheduled BP control, and a
budget-matched control.  The family has five core arms (BP, MF-Joint, MF-Simul,
MF-Seq, and MF-Seq-nom0), two datasets (CIFAR-10 and CIFAR-100), and twenty
paired seeds (42--61), for 200 core runs.  The BP-bestval block adds 40 runs.
The BP-scheduled block uses 50 tuning trials per dataset followed by a
confirmatory comparison.  The Seq-Global block and the budget-matched
Protocol~B block are gated separately; neither is opened by this skeleton.

\subsection{Endpoints and decision rule}
The primary endpoint is the matched-accuracy end-to-end energy ratio at the
pre-registered target.  Secondary endpoints are the per-pass energy ratio, the
pass multiplier, test accuracy at the matched budget, and the stage-wise energy
decomposition.  The family-wise decision uses the pre-registered Holm
procedure over the frozen family.  \wone{insert the exact endpoint definitions
and the family list from PREREG-phase2.md; do not paraphrase the decision rule
in the final manuscript.}

\subsection{Controls and gates}
The BP-bestval control asks whether the result survives a properly tuned BP
baseline.  The BP-scheduled control asks whether a longer, scheduled BP run
closes the gap.  The budget-matched Protocol~B control asks whether MF's
advantage survives when both methods receive the same pass budget derived from
the BP epochs.  Seq-Global is a separate robustness block and remains gated
until its scheduled date.  The gating is part of the pre-registration, not an
analysis choice.

\subsection{Deviations and provenance}
Every deviation is recorded in the dated amendment log before the affected
analysis is read.  The freeze commit, the analysis-code hash, the configuration
manifest hash, and the job-account records are retained with the results.  The
frozen paths are checked against the freeze tag before each block.  No W1
outcome is inspected during health monitoring.

\section{Results}
\label{sec:results}

% Every result in this section is a placeholder.  The table shells in
% analysis/w1_table_shells/ are header-only and contain no W1 outcome.

\subsection{Matched-data-pass energy}
\wone{populate the matched-data-pass energy table from the frozen analysis;
report per-pass energy, accuracy at the matched budget, and the paired-seed
interval.  Do not compare arms on elapsed time alone.}

\subsection{End-to-end energy at the target}
\wone{populate the matched-accuracy energy table and the win-condition figure;
report the measured pair $(a,m)$ with uncertainty, the ratio $R=m/a$, and the
pre-registered decision.}

\subsection{Pass multiplier and convergence}
\wone{report the pass counts to the target, censored runs, and the
convergence curves; separate the schedule effect from the algorithm effect.}

\subsection{Tuned BP and scheduled BP controls}
\wone{report BP-bestval and BP-scheduled results against the same target;
state whether the control changes the sign or the size of the effect.}

\subsection{Budget-matched Protocol B}
\wone{report the Protocol B controls only after the core BP block is complete
and the budget rule has been applied; label every ex-M0 quantity as an
estimate.}

\subsection{Robustness and heterogeneity}
\wone{report seed-level heterogeneity, node/GPU checks, and the pre-registered
sensitivity analyses.  Do not turn health-monitoring elapsed times into an
outcome table.}

\section{Discussion and Limitations}
\label{sec:discussion}

\subsection{What the result would mean}
A measured $R<1$ with a compatible uncertainty interval would support the
end-to-end claim at the tested scale.  A measured $a>1$ with $R\ge1$ would show
that the per-pass advantage is real but insufficient, which is a useful
negative result rather than a failed experiment.  A measured $a\le1$ would
reject the per-pass premise for this implementation and would make the
end-to-end question moot at this scale.

\subsection{Threats to validity}
The main threats are hardware heterogeneity, instrumentation error, tuning
budget mismatch, and the fixed architecture.  The design records node and GPU
identity, retains raw traces, freezes the tuning budget, and pre-registers the
target.  It cannot eliminate node-to-node variation, and it does not vary
width, depth, or device generation.  Those are scope limitations, not
post-hoc excuses.

\subsection{Relation to the conference paper}
The journal version does not replace the conference result; it re-frames it.
The matched-data-pass result answers a per-pass question.  The W1 result
answers the end-to-end question.  Both are reported, with the protocol named
in every table and figure, so that the two claims cannot be conflated.

\section{Conclusion}
\label{sec:conclusion}

Removing the backward pass is a mechanism, not an end-to-end guarantee.  The
cost model in this paper makes the missing step explicit: a forward-only method
wins only when its per-pass advantage exceeds its pass multiplier, after fixed
costs and uncertainty are accounted for.  The W1 family is a pre-registered
test of that condition on a fixed A100 platform.  \wone{insert the one-sentence
W1 conclusion after the frozen analysis; do not preview it here.}

\section*{CRediT authorship contribution statement}
\wone{confirm the final CRediT roles with both authors.}

\section*{Declaration of competing interest}
The authors declare that they have no known competing financial interests or
personal relationships that could have appeared to influence the work reported
in this paper.

\section*{Data and code availability}
The frozen configurations, analysis code, table shells, and provenance records
are available in the project repository.  \wone{insert the repository DOI or
archive identifier and the JPDC data statement after the W1 analysis.}

\section*{Funding}
This work was supported by the National Science Centre (NCN), Poland, under
grant OPUS-29, DEC-2025/57/B/ST6/04377.  \wone{confirm the final funding
statement and any additional institutional support.}

\section*{Acknowledgements}
% Allocation IDs for the JPDC acknowledgement.  Confirm the class-account
% wording before submission; see publication_hygiene.md.
%   - plgoncotherapy2-gpu-a100  (grant plgoncotherapy2, team plggfastsim)
%   - plglscclass26-gpu-a100    (grant plglscclass26; academic/class grant)
%   - PPAM-era allocation PLG/2025/018341
We thank the PLGrid Infrastructure and ACK Cyfronet AGH for computational
resources on the Athena cluster under grant PLG/2025/018341 and allocation
plgoncotherapy2-gpu-a100.  \wone{confirm whether the class-account allocation
plglscclass26-gpu-a100 may be acknowledged; do not include it until the
permitted-use question is resolved.}

\bibliographystyle{elsarticle-harv}
\bibliography{references}

\appendix

\section{Protocol Deviations and Provenance}
\wone{list every dated deviation from PREREG-phase2.md, the freeze tag, the
analysis-code hash, and the configuration-manifest hash.  This appendix must
be generated from the amendment log, not from memory.}

\section{Instrumentation and Measurement Error}
\wone{describe the NVML sampling interval, integration rule, start/stop
boundaries, missing-sample handling, and the stage-sum validation.}

\section{Supplementary Tables}
\wone{place the full per-seed tables and the table-shell schema here; the
header-only shells in analysis/w1\_table\_shells/ are the starting point.}

\end{document}
